% 问题二：问题分析

本题要求对问题 1 检测出的 237 对时频冲突进行消解，使得消解后的用频计划\textbf{不存在任何时频冲突}。

\textbf{消解原则}（题目给定 5 条）：
\begin{enumerate}
    \item 除非特别说明，一般\textbf{不调整使用次数和使用间隔时长}；
    \item 应尽可能\textbf{减少撤销}用频计划；
    \item 尽量\textbf{减少被调整的用频计划数量}；
    \item 按装备类别区分优先级：\textbf{A 类 > B 类 > C 类}，高优先级用频装备的用频计划应尽量保持；
    \item 尽量\textbf{减少调整的幅度}，规定频段、时间调整的最大平移幅度分别为 $10\Delta f$、$5\Delta t$。
\end{enumerate}

\textbf{输入}：150 个原始用频计划 + 237 对冲突关系。

\textbf{输出}：每个计划的处置方案——保留（不变）、调整（频段或时间平移）、或撤销，结果保存到 \texttt{result2.xlsx}。

\textbf{核心约束}：
\begin{enumerate}
    \item 每个计划最多调整 1 个参数（频段或时间），或撤销，不能同时调整多个参数；
    \item 频段平移后须仍在 $[0, 100)$ 范围内，时间起点须 $\geq 0$。
\end{enumerate}

\textbf{优化目标}（分层优先级，对应消解原则）：
\begin{enumerate}
    \item 第 1 层：最小化冲突对数（确保消解后 0 冲突）；
    \item 第 2 层：最小化撤销总数；
    \item 第 3 层：依次最小化 A 类、B 类、C 类撤销数；
    \item 第 4 层：最小化被调整的计划数；
    \item 第 5 层：最小化总调整幅度。
\end{enumerate}
